% ===========================================================================
%  evnx init
% ===========================================================================
\section{evnx init}

\begin{frame}{\cmd{evnx init} — start a project}
  \textbf{Purpose.} Create \file{.env}, \file{.env.example} and a correct
  \file{.gitignore} for a project that has none.

  \vspace{0.6em}
  \begin{block}{Synopsis}
    \texttt{evnx init [OPTIONS]}
  \end{block}

  \vspace{0.4em}
  \begin{tabular}{@{}ll@{}}
    \toprule
    \textbf{Flag} & \textbf{What it does} \\
    \midrule
    \flag{--detect}          & Read the project's \emph{manifests} and scaffold from them \\
    \flag{--from-source}     & Read the project's \emph{code} and scaffold from that \\
    \flag{--with <LIST>}     & Name components directly: \texttt{nextjs,postgresql} \\
    \flag{--blueprint <ID>}  & A named stack, e.g. \texttt{t3\_modern} \\
    \flag{--list-components} & Print everything \flag{--with} accepts, then exit \\
    \flag{-y, --yes}         & Non-interactive; blank unless a stack is named \\
    \flag{--force}           & Replace files that already exist \\
    \flag{--path <PATH>}     & Where to write (default \texttt{.}) \\
    \bottomrule
  \end{tabular}

  \vspace{0.5em}
  {\footnotesize \textbf{Effectively run-once.} A second run refuses rather than
  overwriting a \file{.env.example} you may have hand-maintained — exit \exittwo,
  with \flag{--force} named as the way through.}
\end{frame}

\begin{frame}[fragile]{\cmd{init --detect} — scaffold from the manifest}
  \capture{10-init-detect.txt}
  \realrun{a directory containing only \file{package.json} with \texttt{next}}
\end{frame}

\begin{frame}[fragile]{What \flag{--detect} wrote}
  \capture{10-init-detect-example.txt}

  {\footnotesize Note the \file{.gitignore} line it generated:
  \texttt{.env*, !.env.example, !.env.sample, !.env.template}. The negation is
  the part people get wrong by hand.}
\end{frame}

\begin{frame}[fragile]{\cmd{init --from-source} — scaffold from the code}
  Reads \texttt{process.env.X}, \texttt{os.getenv("X")},
  \texttt{os.environ["X"]} and \texttt{std::env::var("X")}.

  \capture{11-init-from-source.txt}

  {\footnotesize It \emph{refuses} rather than guessing: no terminal and no
  \flag{--yes} is exit \exittwo, and it names the exact command to proceed.}
\end{frame}

\begin{frame}[fragile]{\cmd{init --from-source --yes} — accepted}
  \capture{11b-init-from-source-yes.txt}

  \begin{block}{Why this is the best mode}
    It is derived from what the code \emph{actually reads}, so it answers
    ``is my \file{.env.example} complete?'' from evidence rather than from a
    dependency list. Output is deterministic and each entry carries a
    \texttt{file:line} citation.
  \end{block}
\end{frame}

\begin{frame}[fragile]{\cmd{init --with} — name the stack}
  \capture{12-init-with.txt}
  \realrun{\texttt{evnx init --with nextjs,postgresql,stripe -y}}

  {\footnotesize Discover the vocabulary with \cmd{evnx init --list-components}.}
\end{frame}

\begin{frame}[fragile]{The generated \file{.env}}
  \capturepart{12-init-with-env.txt}{1}{16}

  {\footnotesize \textbf{Changed in v0.5.0:} \file{.env} is written with
  \emph{active assignments}. It previously wrote every line commented as
  \texttt{\# TODO:}, which meant \cmd{evnx validate} — init's own suggested next
  step — reported every variable missing.}
\end{frame}

\begin{frame}[fragile]{When to use which mode}
  \begin{tabular}{@{}lll@{}}
    \toprule
    \textbf{Situation} & \textbf{Use} & \textbf{Why} \\
    \midrule
    Existing codebase      & \flag{--from-source} & Evidence from the code itself \\
    Fresh scaffold         & \flag{--detect}      & Manifest already declares the stack \\
    Known stack, scripted  & \flag{--with}        & Explicit, reproducible, CI-safe \\
    Greenfield, exploring  & (interactive)        & Guided walk \\
    \bottomrule
  \end{tabular}

  \vspace{1em}
  \begin{block}{Re-running is refused, not silently destructive}
    \capture{15-init-rerun.txt}
  \end{block}
\end{frame}
